% 本文件由 LaTeX 排版 Agent 根据有效 translation.json 自动生成。
% author=Apocrypha; work_id=0810; title=般雀·比拉多致提庇留书

\chapter{般雀·比拉多致提庇留书}

% structure: p0001 source=h1 target=omitted reason=page-title-already-rendered-by-parent

\Emph{般雀·比拉多写给罗马皇帝的书信，关于我们的主耶稣基督。}\par

\Person{般雀·比拉多}向\Person{提庇留·凯撒}皇帝请安。\par

关于耶稣基督，我在上一次信中已向您清楚陈明他的案件，最后因百姓的意愿，对他施以残酷的刑罚，我本人多少不情愿且颇为害怕。凭\Person{赫拉克勒斯}起誓，如此虔诚而严谨的人，过去任何时代不曾有，将来也不会有。但这百姓自身的努力，以及所有经师、首领和长老的一致行动，竟要把这真理的使者钉在十字架上，实为惊异。虽然他们自己的先知，以及如同我们那里的西彼拉，曾警告他们不可如此。当他被悬挂时，出现了超自然的征兆，按哲学家们的意见，这些征兆预示着全世界的毁灭。他的门徒正在兴旺，他们的工作与生活规范没有辜负他们的师傅；是的，因他的名而行最有益的事。若不是我担心百姓中即将爆发的叛乱，也许此人至今仍与我们同在。虽然我更多是出于对你尊严的忠诚，而非出于自己的意愿，我没有尽其所能地阻止那无辜的血——它完全摆脱了\Emph{加于它的}所有指控——而不公正地，由于人们的恶意，被出卖并受苦，然而，正如圣经所指示的，这导致他们自身的毁灭。别了。三月二十八日。\par

\SourceRecord{Translated by Alexander Walker.；Ante-Nicene Fathers；Vol. 8.；Edited by Alexander Roberts, James Donaldson, and A. Cleveland Coxe.；Buffalo, NY: Christian Literature Publishing Co.,；1886.；<http://www.newadvent.org/fathers/0810.htm>.}
